% Generated by revision/make_paper_tables.py -- do not edit by hand.
% Regenerate after re-running the Phase 1 task scripts.
\begin{table}[H]
\centering
\begin{threeparttable}
\caption{Absolute error against trivial floors, and behaviour in the high-efficiency
tail where absolute predictions matter for design.}
\label{tab:calibfloor}
\footnotesize
\begin{tabular}{l S[table-format=1.4] S[table-format=1.4] S[table-format=+1.4] c}
\toprule
Predictor & {MAE} & {RMSE} & {Bias} & {95\% CI on MAE} \\
\midrule
constant at training mean & 0.1892 & 0.2103 & +0.1078 & $[0.1851,\,0.1937]$ \\
constant at training median & 0.1227 & 0.1827 & -0.0276 & $[0.1144,\,0.1313]$ \\
constant at held-out mean (oracle floor) & 0.1328 & 0.1806 & +0.0000 & $[0.1263,\,0.1398]$ \\
OptiPrime & 0.0590 & 0.1026 & +0.0127 & $[0.0549,\,0.0629]$ \\
\bfseries PE-RankFormer + isotonic & \bfseries 0.0478 & \bfseries 0.0912 & -0.0038 & $[0.0431,\,0.0523]$ \\
\midrule
\multicolumn{5}{l}{\emph{The calibrated model inside its own top-scoring slices}} \\
\quad top 5\% ($n=1,025$) & \multicolumn{4}{l}{predicts 0.5987, observes 0.5730 --- over-predicts by +0.0257} \\
\quad top 1\% ($n=205$) & \multicolumn{4}{l}{predicts 0.7985, observes 0.7334 --- over-predicts by +0.0651} \\
\bottomrule
\end{tabular}
\begin{tablenotes}[flushleft]\footnotesize
\item Constant baselines are fitted on the 297{,}962 training rows only (median
0.0856, mean 0.2211). The held-out mean (0.1133) is an oracle floor no honest
predictor could use, and is shown only to bound the exercise.
\item The best trivial predictor reaches MAE 0.1227, so the calibrated model's 0.0478
is a 61\% reduction against that floor: the absolute error reflects skill rather
than the scale of a small-valued target.
\item \textbf{The isotonic map over-predicts its own tail.} A user told to expect
0.80 in the model's top percentile should expect 0.73. Isotonic regression has
the least data there and is monotone-constrained, so the final bin absorbs the whole
upper range. Ranking is unaffected \citep{zadrozny2002}; absolute predictions in the tail
should be read as optimistic, which is the regime where a design decision is most likely
to rest on the absolute number \citep{guo2017}.
\item Reproduce with \path{revision/task_1_5_calibration_baselines.py}.
\end{tablenotes}
\end{threeparttable}
\end{table}
